\documentclass[12pt,a4paper]{article}

% Packages for formatting and graphics
\usepackage[utf8]{inputenc}
\usepackage[T1]{fontenc}
\usepackage[english]{babel}
\usepackage{graphicx}
\usepackage{hyperref}
\usepackage{geometry}
\usepackage{booktabs}
\usepackage{longtable}
\usepackage{listings}
\usepackage{xcolor}

% Page geometry
\geometry{margin=1in}

% Code listing style
\lstset{
    basicstyle=\ttfamily\small,
    breaklines=true,
    frame=single,
    backgroundcolor=\color{gray!5},
    keywordstyle=\color{blue},
    commentstyle=\color{green!50!black},
    stringstyle=\color{red},
    captionpos=b
}

\title{Assignment 6: Making the Repository Great \\ \large Poetry: A Conversational Recommender System}
\author{Zakhar (Solo Developer) \\ Innopolis University \\ \href{mailto:z.user@innopolis.university}{z.user@innopolis.university} \\ GitHub: \href{https://github.com/Potter32111}{Potter32111}}
\date{\today}

\begin{document}

\maketitle

\section*{Who Did What During This Sprint}
\begin{itemize}
    \item \textbf{Zakhar}: Ported all documentation to specialized Markdown files in the \texttt{docs/} directory. Rewrote the root \texttt{README.md} for better clarity/navigation. Generated project logo. Created \texttt{CHANGELOG.md} and applied MIT License. Finalized MVP v2.5 with voice recitation improvements.
\end{itemize}

\newpage
\tableofcontents
\newpage

\section{Introduction}
This report summarizes the efforts to finalize the project's repository structure, enhance user-facing documentation, and ensure a smooth transition of the \textbf{Poetry Recommender} system to the customer. The focus of this sprint was on "Making the Repo Great" by adhering to professional documentation standards and providing a clear roadmap for future development.

\section{MVP Roadmap}
The project has evolved through several Minimum Viable Product (MVP) stages. The roadmap is now documented and linked within the main repository.

\textbf{Link to Roadmap}: \href{https://github.com/Potter32111/poetry-recommender#feature-roadmap}{GitHub Feature Roadmap}

\subsection{Current Status (MVP v2.5)}
Currently, the system supports:
\begin{itemize}
    \item Personalized English/Russian poem recommendations.
    \item SM-2 based spaced repetition scheduling.
    \item Offline Voice Recitation check (Vosk STT).
    \item Comprehensive internal documentation for all system components.
\end{itemize}

\section{Sprint Tracking}
During this sprint (Sprint 6), the following Story Points (SP) were addressed:

\begin{longtable}{@{}p{0.7\textwidth}p{0.2\textwidth}p{0.1\textwidth}@{}}
\toprule
\textbf{Task / Issue} & \textbf{Status} & \textbf{SP} \\ \midrule
Create Documentation Structure (\texttt{docs/}) & Completed & 5 \\
Rewrite README with Logo \& Metadata & Completed & 3 \\
Update Architecture Diagrams (Mermaid) & Completed & 3 \\
Implement CHANGELOG and MIT License & Completed & 2 \\
Bugfix: Russian Voice Check Accuracy & Completed & 5 \\
Finalize Assignment 6 Report & In Progress & 2 \\
\midrule
\textbf{Total Completed} & & \textbf{18} \\
\bottomrule
\end{longtable}

\section{ReadMe Sections Overview}
The \texttt{README.md} was completely restructured to serve as the single source of truth for the project. Key sections include:

\begin{itemize}
    \item \textbf{Header}: Features a high-quality logo, one-liner, and direct links to the deployed bot and demo video.
    \item \textbf{Context Diagram}: A Mermaid-based visualization of system stakeholders (User, Telegram, Backend, Database).
    \item \textbf{Feature Roadmap}: A checklist tracking past achievements (MVP v1-v2.5) and future goals (MVP v3+).
    \item \textbf{Usage \& Deployment}: Clear, tested instructions for reproducibility.
    \item \textbf{Documentation Index}: Hyperlinks to specialized Markdown files for Development, QA, and Architecture.
\end{itemize}

\section{Conclusion \& Next Steps}
The repository is now in a "world-class" state, with all documentation easily accessible and a clear path forward defined in the roadmap. The next sprint will focus on MVP v3, including semantic search capabilities and further gamification of the learning experience.

\end{document}
